\documentclass[11pt,a4paper]{report}

\usepackage[utf8]{inputenc}
\usepackage[T1]{fontenc}
\usepackage[english]{babel}
\usepackage{amsmath}
\usepackage{graphicx}
\usepackage{siunitx}
\usepackage{booktabs}

% --- Literaturverzeichnis ---
\usepackage[backend=biber, style=numeric, sorting=none]{biblatex}
\addbibresource{references.bib}

\title{LaTeX Toolchain Test}
\author{Konrad Braun}
\date{\today}

\begin{document}

\maketitle

\chapter{Test}

\section{Text, units and citations}

The hydration enthalpy of zinc is approximately
\SI{-2046}{\kilo\joule\per\mole}, compared with roughly
\SI{-520}{\kilo\joule\per\mole} for lithium. The measured ionic
conductivity of the ND@GF separator is \SI{2.94}{\siemens\per\metre}.

Dendrite formation remains the central failure mode of the zinc metal
anode~\cite{yu_ten_2023}. Separator engineering has emerged as one
lever to address it~\cite{gu_role_2025}, and the underlying
nucleation mechanism has been studied in
detail~\cite{zhang_unveiling_2024}.

Multiple sources can be cited at once~\cite{yu_ten_2023,gu_role_2025}.
Textual citations also work: \textcite{zhang_unveiling_2024} investigated
dendrite nucleation directly.

\section{Equation}

Sand's time relates the anion transference number to the onset of
dendrite growth:

\begin{equation}
    \tau_\text{Sand} = \pi D \left( \frac{z_c c_0 F}{2 J t_a} \right)^{2}
    \label{eq:sand}
\end{equation}

Equation~\eqref{eq:sand} shows that $\tau_\text{Sand}$ scales with
$1/t_a^{2}$.

\section{Table}

\begin{table}[h]
    \centering
    \caption{Measured separator properties.}
    \label{tab:properties}
    \begin{tabular}{lcc}
        \toprule
        Property & GF & ND@GF \\
        \midrule
        Ionic conductivity (\si{\siemens\per\metre}) & 2.83 & 2.94 \\
        Transference number $t_+$ & 0.51 & 0.69 \\
        Activation energy (\si{\kilo\joule\per\mole}) & 30.8 & 19.3 \\
        \bottomrule
    \end{tabular}
\end{table}

\printbibliography

\end{document}
